\documentclass[10pt,a4paper]{amsart}
\usepackage{amsmath,amssymb,amsthm,graphicx}
% Retained preamble for the Stacks Project chapter snapshot in this directory.
% The chapter's own preamble.tex is a copy of the shared Stacks Project preamble; the block
% below is sliced MECHANICALLY, line for line, from it. Nothing was added, removed or
% redefined; the wrapper class, the xr-hyper/hyperref wiring and the externaldocument list
% are omitted because the assembler supplies the document class.
%
% lines taken from preamble.tex: 3, 5, 9, 11, 12, 156, 157, 158, ... (51 lines total)
%
% package coverage restored from the chapter's own preamble (2026-10-13): xy, xy-2cell
% package coverage restored from the chapter's own preamble (2026-10-14): verbatim
%   the chapter's preamble.tex loads \usepackage{verbatim} on line 8 and then defines
%   reference/slogan/history as \comment environments; without the package \comment is
%   undefined and any card whose counted spans use one of the three environments fails
%   with "Undefined control sequence" (L44184, lemma-representable-over-space).

\usepackage{verbatim}
\usepackage[all]{xy}
\xyoption{2cell}
\UseAllTwocells

\newenvironment{reference}{\comment}{\endcomment}
\newenvironment{slogan}{\comment}{\endcomment}
\newenvironment{history}{\comment}{\endcomment}
\theoremstyle{plain}
\newtheorem{theorem}[subsection]{Theorem}
\newtheorem{proposition}[subsection]{Proposition}
\newtheorem{lemma}[subsection]{Lemma}
\theoremstyle{definition}
\newtheorem{definition}[subsection]{Definition}
\newtheorem{example}[subsection]{Example}
\newtheorem{exercise}[subsection]{Exercise}
\newtheorem{situation}[subsection]{Situation}
\theoremstyle{remark}
\newtheorem{remark}[subsection]{Remark}
\newtheorem{remarks}[subsection]{Remarks}
\numberwithin{equation}{subsection}
\def\lim{\mathop{\mathrm{lim}}\nolimits}
\def\colim{\mathop{\mathrm{colim}}\nolimits}
\def\Spec{\mathop{\mathrm{Spec}}}
\def\Hom{\mathop{\mathrm{Hom}}\nolimits}
\def\Ext{\mathop{\mathrm{Ext}}\nolimits}
\def\SheafHom{\mathop{\mathcal{H}\!\mathit{om}}\nolimits}
\def\SheafExt{\mathop{\mathcal{E}\!\mathit{xt}}\nolimits}
\def\Sch{\mathit{Sch}}
\def\Mor{\mathop{\mathrm{Mor}}\nolimits}
\def\Ob{\mathop{\mathrm{Ob}}\nolimits}
\def\Sh{\mathop{\mathit{Sh}}\nolimits}
\def\NL{\mathop{N\!L}\nolimits}
\def\CH{\mathop{\mathrm{CH}}\nolimits}
\def\proetale{{pro\text{-}\acute{e}tale}}
\def\etale{{\acute{e}tale}}
\def\QCoh{\mathit{QCoh}}
\def\Ker{\mathop{\mathrm{Ker}}}
\def\Im{\mathop{\mathrm{Im}}}
\def\Coker{\mathop{\mathrm{Coker}}}
\def\Coim{\mathop{\mathrm{Coim}}}
\DeclareMathSymbol{\boxtimes}{\mathbin}{AMSa}{"02}
\def\QCohstack{\mathcal{QC}\!\mathit{oh}}
\def\Cohstack{\mathcal{C}\!\mathit{oh}}
\def\Spacesstack{\mathcal{S}\!\mathit{paces}}
\def\Quotfunctor{\mathrm{Quot}}
\def\Hilbfunctor{\mathrm{Hilb}}
\def\Curvesstack{\mathcal{C}\!\mathit{urves}}
\usepackage{cleveref}
\crefname{originallemma}{original lemma}{original lemmas}
\Crefname{originallemma}{Original lemma}{Original lemmas}
\makeatletter
\newlabel{lemma-properties-diagonal}{{02X4}{source}{}{sourcepointer-1}{}}
\newlabel{lemma-properties-diagonal@cref}{{[originallemma][0][]02X4}{[1][0][]source}{}{}{}}
\makeatother
\crefname{originallemma}{original lemma}{original lemmas}
\Crefname{originallemma}{Original lemma}{Original lemmas}
\makeatletter
\newlabel{lemma-quotient}{{02Z2}{source}{}{sourcepointer-2}{}}
\newlabel{lemma-quotient@cref}{{[originallemma][0][]02Z2}{[1][0][]source}{}{}{}}
\makeatother
\title{Quotients by a finite group action inherit quasi-separatedness and separatedness}
\author{The Stacks Project Authors}
\date{stacks-project:spaces.tex:local-snapshot-20261012; GNU FDL 1.2 or later}
\begin{document}
\maketitle
\setcounter{subsection}{1}
\noindent\textit{Editorial note.} All delivered mathematical text is the author's own, copied from the pinned Stacks Project chapter file spaces.tex; no statement or proof step was rewritten, reordered or removed.\par\medskip
\noindent\textit{Editorial note.} The statement inherits the notation of the preceding quotient lemma through a reference; the reference is recorded as a non-load-bearing navigation pointer while the counted statement stays self-contained in wording.\par\medskip
\noindent\textit{Editorial note.} The two references of the counted spans are recorded as non-load-bearing navigation pointers; the argument that consumes them stays inside the counted proof span.\par\medskip
\noindent\textit{Editorial note.} Preamble: sources/stacks-spaces/retained-preamble.tex, sliced mechanically from the chapter's own preamble.tex.\par\medskip
\noindent\textit{Allowed standard prerequisites (editorial).}\begin{itemize}
\item The quotient by a group action is an etale equivalence relation. Lemma lemma-quotient (2463-2487, produced as L44177), used at proof lines 2549-2550. Supplies the notation and assumptions.
\item Reduction of diagonal properties of a quotient to the map j. Lemma lemma-properties-diagonal (2296-2308, produced as L44181), used at proof lines 2559-2561. Reduces the claims to properties of j.
\end{itemize}\medskip
\ifdefined\hypertarget\hypertarget{sourcepointer-1}{}\fi
\noindent\textit{Source reference.} Original lemma 02X4: focused-seeds/AGP/sources/stacks-snapshot-20261011/tags.txt, source lines 3827--3827. Navigation pointer into the Stacks Project; the standard fact it supplies is applied inside the counted proof span and the argument resting on it stays there.\par\medskip
\ifdefined\hypertarget\hypertarget{sourcepointer-2}{}\fi
\noindent\textit{Source reference.} Original lemma 02Z2: focused-seeds/AGP/sources/stacks-snapshot-20261011/tags.txt, source lines 3897--3897. Navigation pointer into the Stacks Project; the standard fact it supplies is applied inside the counted proof span and the argument resting on it stays there.\par\medskip
\begin{lemma}
\label{lemma-quotient-finite-separated}
Notation and assumptions as in Lemma \ref{lemma-quotient}.
Assume $G$ is finite. Then
\begin{enumerate}
\item if $U \to S$ is quasi-separated, then $U/G$ is quasi-separated
over $S$, and
\item if $U \to S$ is separated, then $U/G$ is separated over $S$.
\end{enumerate}
\end{lemma}

\begin{proof}
In the proof of Lemma \ref{lemma-properties-diagonal}
we saw that it suffices to prove the
corresponding properties for the morphism $j : R \to U \times_S U$.
If $U \to S$ is quasi-separated, then for every affine open $V \subset U$
which maps into an affine of $S$
the opens $g(V) \cap V$ are quasi-compact. It follows that $j$ is
quasi-compact.
If $U \to S$ is separated, the diagonal $\Delta_{U/S}$ is a closed
immersion. Hence $j : R \to U \times_S U$ is a finite coproduct
of closed immersions with disjoint images. Hence $j$ is a closed immersion.
\end{proof}

\end{document}
